% -----------------------------------------------------------------------
\section{Diseño Experimental}
\label{sec:metodologia}

Esta sección describe el diseño experimental del estudio. Primero se
presenta la organización del repositorio y su arquitectura de software;
luego, las instancias de prueba, las heurísticas constructivas que generan
las soluciones iniciales y la métrica con que se compara el rendimiento
de los tres niveles de TS.

\subsection{Repositorio y estructura de código}

El repositorio está organizado en las siguientes carpetas principales:

\begin{itemize}
  \item \texttt{instances/} — instancias TSPLIB en formato estándar
    (extensiones \texttt{.kp}, \texttt{.tsp}, \texttt{.vrp}, \texttt{.fsp},
    \texttt{.jss}).
  \item \texttt{tabu\_search/\{problema\}/} — quince módulos Python,
    uno por combinación (problema, nivel), con interfaz uniforme
    \texttt{solve(instance\_data, \ldots) → dict}.
  \item \texttt{src/} — orquestador (\texttt{main.py}), cargador de
    instancias (\texttt{loader.py}), calculadora de métricas
    (\texttt{metrics.py}), generador de informes (\texttt{reporter.py})
    y utilidades (\texttt{utils/}).
  \item \texttt{models/} — modelos exactos de referencia: formulaciones en
    AMPL y solvers en Python (programación dinámica para Knapsack y
    Held-Karp~\cite{held1962dynamic} para TSP).
  \item \texttt{output/\{problema\}/} — resultados de cada experimento.
  \item \texttt{notebooks/} — cuadernos Jupyter para análisis y figuras.
  \item \texttt{paper/} — este documento y sus recursos.
\end{itemize}

La Figura~\ref{fig:arquitectura_repo} esquematiza la arquitectura y el
flujo de datos entre estos componentes.

\begin{figure}[t]
  \centering
  \begin{tikzpicture}[
      node distance=0.5cm and 0.95cm,
      comp/.style={draw, rounded corners=2pt, align=center,
                   font=\scriptsize\ttfamily, minimum width=2.15cm,
                   minimum height=0.85cm, fill=gray!8},
      etapa/.style={font=\scriptsize\itshape, text=gray!90!black},
      fl/.style={->, >=stealth, thick}]
    % Fila 1: entrada y núcleo
    \node[comp] (inst) {instances/\\{\tiny .kp .tsp .vrp .fsp .jss}};
    \node[comp, right=of inst] (src) {src/\\{\tiny main.py, loader.py}};
    \node[comp, right=of src] (ts) {tabu\_search/\\{\tiny 15 módulos solve()}};
    % Fila 2: salida y análisis
    \node[comp, below=of ts] (out) {output/\\{\tiny resultado.json}};
    \node[comp, left=of out] (rep) {src/\\{\tiny metrics.py, reporter.py}};
    \node[comp, left=of rep] (nb) {notebooks/ y\\paper/\\{\tiny tablas y figuras}};
    % Flujos
    \draw[fl] (inst) -- node[above, etapa]{carga} (src);
    \draw[fl] (src) -- node[above, etapa]{ejecuta} (ts);
    \draw[fl] (ts) -- node[right, etapa]{guarda} (out);
    \draw[fl] (out) -- node[above, etapa]{agrega} (rep);
    \draw[fl] (rep) -- node[above, etapa]{genera} (nb);
  \end{tikzpicture}
  \caption{Arquitectura del repositorio y flujo de datos: el orquestador
           carga las instancias, ejecuta los quince módulos de TS y
           persiste los resultados, desde los cuales se generan métricas,
           tablas y figuras.}
  \label{fig:arquitectura_repo}
\end{figure}

Todos los módulos TS exponen la misma interfaz de devolución:
\texttt{\{value, solution, iterations, time\_s, history\}}, interfaz uniforme
que permite al orquestador tratar los cinco problemas de forma homogénea.
La Figura~\ref{fig:uml_clases} presenta el diagrama de clases UML de los
módulos: \texttt{runner.py} depende únicamente de la interfaz
\texttt{solve()}, que los quince módulos de TS realizan y que se carga
mediante importación dinámica.

\begin{figure}[t]
  \centering
  \begin{tikzpicture}[
      clase/.style={draw, rectangle split, rectangle split parts=2,
                    align=left, font=\scriptsize\ttfamily, fill=gray!8,
                    inner sep=2.5pt, text width=3.75cm, anchor=north},
      dep/.style={->, >=stealth, dashed},
      real/.style={-{Triangle[open, length=5pt, width=6pt]}, dashed},
      rot/.style={font=\tiny\itshape, text=gray!90!black}]
    % Columna izquierda: orquestación y algoritmos
    \node[clase] (main) at (0,0)
      {\centering «módulo»\\\textbf{src/main.py}
       \nodepart{two}+ main()\\+ run\_all(seeds)};
    \node[clase] (run) at (0,-1.75)
      {\centering «módulo»\\\textbf{src/runner.py}
       \nodepart{two}+ run\_experiment(\\\ \ problema, instancia,\\\ \ algoritmo, seed): dict};
    \node[clase] (ifz) at (0,-3.9)
      {\centering «interface»\\\textbf{SolverTS}
       \nodepart{two}+ solve(instance\_data,\\\ \ ..., seed): dict};
    \node[clase] (impl) at (0,-5.75)
      {\centering «módulo»\\\textbf{tabu\_search/\{problema\}/}
       \nodepart{two}ts\_basico.py\\ts\_intermedio.py\\ts\_avanzado.py};
    % Columna derecha: servicios
    \node[clase] (rep) at (4.4,0)
      {\centering «módulo»\\\textbf{src/reporter.py}
       \nodepart{two}+ generate\_report()\\\mbox{}};
    \node[clase] (load) at (4.4,-1.75)
      {\centering «módulo»\\\textbf{src/loader.py}
       \nodepart{two}+ load\_instance(ruta):\\\ \ dict\\\mbox{}};
    \node[clase] (met) at (4.4,-3.9)
      {\centering «módulo»\\\textbf{src/metrics.py}
       \nodepart{two}+ gap(valor, bks, min)\\+ summarize(...): dict};
    % Dependencias y realización
    \draw[dep] (main) -- (rep);
    \draw[dep] (main) -- (run);
    \draw[dep] (run) -- (load);
    \draw[dep] (run) -- node[right, rot]{importación dinámica} (ifz);
    \draw[dep] (run.south east) -- (met.north west);
    \draw[real] (impl) -- (ifz);
  \end{tikzpicture}
  \caption{Diagrama de clases UML de los módulos del repositorio. Las
           flechas punteadas indican dependencia; la flecha con triángulo
           hueco, la realización de la interfaz \texttt{solve()} por los
           quince módulos de TS.}
  \label{fig:uml_clases}
\end{figure}

\subsubsection*{Estructura de la carpeta de salida}

Cada ejecución genera un directorio propio dentro de
\texttt{output/\{problema\}/} con el patrón
\texttt{\{instancia\}\_\{algoritmo\}\_\{fecha\}/}, que contiene
\texttt{resultado.json} (métricas: valor, GAP, iteraciones, tiempo e
historial de convergencia) y \texttt{solucion.txt} (la mejor solución
encontrada en formato legible).
En la raíz de \texttt{output/} se generan además las tablas agregadas
(\texttt{tabla\_\{problema\}.csv}, \texttt{tabla\_consolidada.csv}) y el
script \texttt{regenerar\_figuras.py} que reconstruye las figuras del
artículo:

\begin{figure}[H]
\footnotesize
\begin{verbatim}
output/
|- tsp/
|  |- kroA100_ts_avanzado_<fecha>/
|  |  |- resultado.json
|  |  `- solucion.txt
|  `- ...
|- knapsack/  vrp/  flowshop/  jobshop/
|- tabla_tsp.csv ... tabla_consolidada.csv
`- regenerar_figuras.py
\end{verbatim}
\vspace{-0.5em}
\caption{Estructura de la carpeta de salida \texttt{output/}.}
\label{fig:output_tree}
\end{figure}

\subsection{Instancias de prueba}

La Tabla~\ref{tab:instancias} lista las 24 instancias empleadas.
Se eligieron instancias de tamaño moderado de TSPLIB para mantener tiempos
de ejecución razonables —del orden de uno a dos minutos por experimento—
sin perder representatividad.

\begin{table}[tb]
\caption{Instancias de prueba}
\label{tab:instancias}
\centering
\footnotesize
\setlength{\tabcolsep}{4pt}
\begin{tabular}{llrrl}
\toprule
\textbf{Problema} & \textbf{Instancia} & \textbf{Tamaño} &
  \textbf{BKS} & \textbf{Fuente} \\
\midrule
\multirow{5}{*}{Knapsack}
  & knapPI\_11 & 100 ítem. & 3950  & Pisinger~\cite{pisinger2005hard} \\
  & knapPI\_12 & 100 ítem. & 3752  & \\
  & knapPI\_13 & 100 ítem. & 3468  & \\
  & knapPI\_14 & 100 ítem. & 3856  & \\
  & knapPI\_15 & 100 ítem. & 3732  & \\
\midrule
\multirow{5}{*}{TSP}
  & eil51      &  51 ciud. &  426  & TSPLIB~\cite{reinelt1991tsplib} \\
  & berlin52   &  52 ciud. & 7542  & \\
  & eil76      &  76 ciud. &  538  & \\
  & kroA100    & 100 ciud. & 21282 & \\
  & kroB100    & 100 ciud. & 22141 & \\
\midrule
\multirow{4}{*}{CVRP}
  & vrpnc1     &  50 clien.&  524  & Christofides~\cite{christofides1979vehicle} \\
  & vrpnc2     &  75 clien.&  835  & \\
  & vrpnc3     & 100 clien.&  826  & \\
  & vrpnc4     & 150 clien.& 1028  & \\
\midrule
\multirow{5}{*}{Flowshop}
  & ta001      & 20$\times$5  & 1278  & Taillard~\cite{taillard1993benchmarks} \\
  & ta011      & 20$\times$10 & 1582  & \\
  & ta021      & 20$\times$20 & 2297  & \\
  & ta031      & 50$\times$5  & 2724  & \\
  & ta041      & 50$\times$10 & 2991  & \\
\midrule
\multirow{5}{*}{Job Shop}
  & ft10       & 10$\times$10 &  930  & Fisher\&Thompson~\cite{fisher1963probabilistic} \\
  & la01       &  10$\times$5 &  666  & Lawrence~\cite{lawrence1984resource} \\
  & la06       &  15$\times$5 &  926  & \\
  & la11       &  20$\times$5 & 1222  & \\
  & la21       &  15$\times$10& 1046  & \\
\bottomrule
\end{tabular}
\end{table}

\subsection{Soluciones iniciales}

Cada problema emplea una heurística constructiva para generar la solución
inicial:

\begin{itemize}
  \item \textbf{Knapsack}: greedy por ratio ganancia/peso decreciente.
  \item \textbf{TSP}: vecino más cercano desde el nodo 0.
  \item \textbf{CVRP}: ahorros de Clarke-Wright~\cite{clarke1964scheduling}.
  \item \textbf{Flowshop}: heurística NEH~\cite{nawaz1983heuristic}.
  \item \textbf{Job Shop}: decodificador de Giffler-Thompson con prioridad
    por tiempo total de procesamiento decreciente.
\end{itemize}

\subsection{Métricas}

El rendimiento se mide con el \emph{GAP} respecto al mejor valor conocido
(BKS):

\begin{equation}
  \text{GAP}(\%) =
  \begin{cases}
    \dfrac{f^* - \text{BKS}}{\text{BKS}} \times 100
      & \text{minimización} \\[6pt]
    \dfrac{\text{BKS} - f^*}{\text{BKS}} \times 100
      & \text{maximización}
  \end{cases}
  \label{eq:gap}
\end{equation}

donde $f^*$ es el mejor valor encontrado por el algoritmo.
Un GAP igual a 0\,\% significa que el algoritmo alcanzó exactamente el BKS;
un GAP positivo mide cuánto peor fue la solución obtenida respecto al BKS.
Un GAP negativo indicaría una solución mejor que el BKS registrado; no se
observa en ninguno de los 360 experimentos.
Cada experimento se ejecuta con 5 semillas independientes
($s \in \{42, 123, 456, 789, 1337\}$); se reportan media y desviación
estándar.
Aunque cada corrida de TS es determinista una vez fijada la semilla, la
semilla determina los componentes estocásticos del algoritmo (desempates,
perturbaciones de reinicio y diversificación), de modo que cada semilla
genera una corrida independiente.
Se seleccionan 5 semillas —y no un número mayor— como equilibrio entre costo
computacional (360 experimentos totales) y estimación de la variabilidad;
la baja desviación estándar observada respalda esta decisión.
